
\chapter{Continuous Improvement}
\label{sec:continuousimprovement}

%   Savas
Within the bank's Shared Service Center IT operating model, Continous Improvement represents a critical success factor focused on the systematic and sustained enhancement of IT service performance through structured, data-driven decision-making. It establishes improvement as an embedded capability rather than a set of isolated initiatives, ensuring that IT services, processes and supporting technologies evolve in alignment with business objectives previously defined \ref{sec:goals}. \\
In this context, Continual Improvement is not limited to reactive problem resolution but extends to proactive optimisation of service performance, cost structures, and delivery speed across all IT service management domains.

This \gls{csf} is monitored through the consistent use of performance metrics, service reporting and analytical feedback loops, that convert operational data into actionable improvement opportunities. In a complex, geographically distributed banking environment with multiple acquired entities, Continual Improvement, becomes essential to reduce performance variability, eliminate systemic inefficiencies and ensure that IT services remain scalable and compliant with different regulations. 
Moreover, it also directly supports the bank's strategic ambition to integrate heterogeneous IT landscapes by promoting standardisation practices measurements and reinforcing a unified performance management culture across all divisions.

A key characteristic of this CSF is its dependency on mature measurement and governance capabilities. Improvement decisions are expected to be grounded in quantifiable evidence derived from ITSM processes rather than subjective assessment. As a result, Continual Improvement is tied with the organisation's ability to capture reliable service data, analyse trends and then, to implement corrective actions through controlled change mechanisms. Without this closed-loop structure, improvement efforts risks to become inconsistent and misaligned with organization goals.


\section{Service Costs}
\label{subs:servicecosts}
%   Savas


Represents the organization's ability to deliver IT services in a cost-effective and efficient way. Focusing on optimization of several resources such as financial, technological and human ones, while maintaining the required service quality. \\
An effective and proactive approach oriented to cost management helps the IT Department to ensure that services provide the maximum value to the business withouth unnecessary expenses. This is a core pillar for supporting the bank's expansion and IT Department as ``Shared Service Center'' strategy.
Achieving this CSF requires continuous monitoring, resource optimization and informed financial decision-making. In the end, controlling service costs contributes to improve profitability and long-term business sustainability.



In the context of the case study, this factor is particularly relevant due to the bank's international expansion, in terms of acquisitions of others institutions but also as collaborations with partners and customers. By measuring the total cost of individual services, the bank can assess whether service delivery remains economically sustainable while respecting the performance and quality expected by stakeholders.



\subsection{IT Service costs distribution}
\label{sub:itservicecost}

Overral, Service Costs is the objective is to balance operational efficiency with service performance and customer expectations. Where ``\textbf{Total IT Service Cost per Service}'' (TISCS) works as KPI to support the Continuous Improvement CSF~\ref{sec:continuousimprovement}.

This metric measures the total financial cost associated with delivering, operating, supporting and maintaining an individual IT service over a defined reporting period. \\
Within the calcualtion are includes both direct costs, such as infrastructure, software licensing, personnel, supplier expenditures; than indirect costs, such as shared facilities, administrative overhead and common support services. \\

Unlike traditional budgetary measures, that provide only an aggregated view of expenditure, this metric enables service-level cost transparency and supports informed managerial decision-making.
Furthermore, it allows the organization to identify cost inefficiencies, compare service performance across business units and international operations, evaluate the financial impact of improvement initiatives; in the end allows to optimize resource allocation. \\


Several IT Service Management processes directly influence this indicator, such as:
\begin{itemize}
    \item \textbf{Service Portfolio Management}: affects the metric by determining which services are offered, maintained, or retired, thereby influencing the overall service cost structure.
    \item \textbf{Financial Management}: for IT Services plays a critical role in budgeting, accounting, cost allocation, and financial reporting associated with individual services.
    \item \textbf{Service Level Management}: impacts service costs through the definition of service quality targets and performance requirements, which determine the resources necessary for service delivery.
    \item \textbf{Change, Release and Deploy Management}: influence the metric by affecting the cost of implementing new or modified services and introducing operational efficiencies or inefficiencies. 
    % Service Asset and Configuration Management contributes by providing accurate information regarding infrastructure assets and their associated costs.  -- not too many i guess, such ash
    \item \textbf{Continual Service Improvement}: directly affects this metric through the identification and implementation of initiatives aimed at reducing operational costs, improving efficiency, and increasing the overall value delivered by IT services.

\end{itemize}
   

In the end, this metric provides a comprehensive and strategically aligned measure of cost effectiveness, supporting both the Bank's Continuous Improvement objectives than its broader expansion.

    \begin{table}[h]
        \centering
        \renewcommand{\arraystretch}{1.6}
        \setlength{\arrayrulewidth}{0.6pt}

        \begin{tabularx}{\textwidth}{|X|X|}
        \hline
        \multicolumn{2}{|c|}{\cellcolor{mcblue}\textcolor{white}{\textbf{Total IT Service Cost per Service}}} \\ \hline

        \rowcolor{mclightgray}
        \textbf{Description} & \textbf{Unit} \\ \hline
        Measures the total cost incurred by the IT organization to deliver, operate, support, and maintain an IT service during a defined reporting period. &
        Currency per service, per reporting period \\ \hline

        \rowcolor{mclightgray}
        \textbf{Formula} & \textbf{Frequency} \\ \hline
        Direct Service Costs + Allocated Indirect Costs &
        Monthly/Quarterly/Annually. \\ \hline

        \end{tabularx}

        \caption{Service costs - TISCS}
        \label{tab:measurement_ITcostperservice}
    \end{table}

In the formulation of this KPI metric, Direct Service Costs and Allocated Indirect Cost elements can be listed as following:

\textbf{Direct Service Costs include:}
\begin{itemize}
    \item Infrastructure and hardware costs
    \item Software licensing costs
    \item Application maintenance costs
    \item Personnel and support staff costs
    \item Third-party supplier and outsourcing costs
\end{itemize}

\textbf{Allocated Indirect Costs include:}
\begin{itemize}
    \item Shared infrastructure costs
    \item Service management tooling costs
    \item Facilities and data center costs
    \item Administrative overhead costs
    \item Shared services costs
\end{itemize}


This metric measures the total cost incurred by the IT organization to deliver, operate, support, and maintain a specific IT service during a defined reporting period. The calculation includes all direct costs attributable to the service as well as allocated indirect costs arising from shared resources and support functions.

%   DONE but check again



\section{Backlog}
\label{subs:backlog}

%   Alberto



Backlog represent a structured and prioritized repository of improvements initiatives taken or identified through service performance measurement methods like \gls{kpi} monitoring, incidents, problem analysis or customer feedback and reviews.
Provide a governance mechanism that enables the organization to systematically collect, evaluate, pripritize and monitor the improvement process and opportunities in a shared system rather than operate as an organizational silos.
In fact, backlog represent a key tool for the transformation of the IT Department into a ``Shared Service Center'', because improvement initiatives are continuously generated from the monitoring of KPI across various processes. The adoption of a welformed backlog gives the IT the possibility to provide and suggest different initiatives thus to ensure the investments of resource and operation that contribute directly to the organization's goals and  support informed decision-making.

Initiatives are prioritized by different factors, such as:
\begin{itemize}
    \item \textbf{Business value} that provide to the customers
    \item \textbf{Strategic alignment} toward organisation's goals
    \item \textbf{Implementation effort}
    \item \textbf{Operational risk}
    \item \textbf{Expected benefits}
    \item \textbf{Resouce ability}
\end{itemize}

The backlog must be continuously reviewed as part of the \textit{Continual Service Improvement} process, allowing new opportunities to be integrated while completed activities can be evaluated against predefined performance indicators.
The outcome of each implemented improvement is measured through opportune KPIs to verify and then support the validation of the achievement of the expected benefits. Moreover, can serve to identify additional optimization and creating a feedback loop that promotes organizational learning and maturity.


Applied to the Bank current ITSM environment, the adoption of an improvement backlog brings several benefits to the services provided, such as:

\begin{itemize}
    \item Enrich and consolidate a Knowledge Management
    \item Increase customer involvement during Service Transition
    \item Introduce formal change evaluation process 
    \item Promote proactive Problem Management
    \item Standardize operational process across international branches and improve collaboration between departments
\end{itemize}

Together, these initiatives contribute the reducing of operational waste while improving service quality, increasing process efficiency and supporting Banks's international growth strategy through a culture of continuous improvement.


\subsection{Priority distribution}
\label{sub:prioritydistribution}
Backlog priority distribution measures the compositiion of open tickets, grouping them by their priority. This metric can be applied to different subjects, such as problems or change requests.

With this KPI~\gls{kpi}, the IT can monitor a rank of the amount of task for each division in way to plan and coordinate different resources in order to accomplish first critical tasks and to direct others requests.

It affects several \gls{itsm} process for each division, for instance:

\begin{itemize}
    \item \textbf{Incident Management} for incidents 
    \item \textbf{Problem Management} for problems and known errors
    \item \textbf{Change Enablement} for change requests open or request awaiting to be accomplished
\end{itemize}

\begin{table}[H]
    \centering
    \renewcommand{\arraystretch}{1.6}
    \setlength{\arrayrulewidth}{0.6pt}

    \begin{tabularx}{\textwidth}{|X|X|}
    \hline
    \multicolumn{2}{|c|}{\cellcolor{mcblue}\textcolor{white}{\textbf{Backlog Priority Distribution}}} \\ \hline

    \rowcolor{mclightgray}
    \textbf{Description} & \textbf{Unit} \\ \hline
    Helps determine whether unresolved work is concentrated on high-impact issues and supports prioritization decisions. The organization should aims to lower percentage on the priority level. &
    Percentage. \\ \hline

    \rowcolor{mclightgray}
    \textbf{Formula} & \textbf{Frequency} \\ \hline
    \[
     \frac{\text{Open tickets of priority P}}{Total open tickets} \times 100
    \]
    
    \textit{calculated for each priority level}
    
     &
    Weekly (operational team) \& Montly (as management reporting). \\ \hline

    \end{tabularx}

    \caption{Measurement - Backlog Priority Distribution}
    \label{tab:measurement_backlogPriorityDistribution}
\end{table}



\subsection{Task completed from last report}
\label{sub:taskcompleted}
Another metric, combined with the before Priority Distribution \ref{sub:prioritydistribution}, provide the work tasks successfully completed during the reporting period. \\
This provide the performance of the various department toward the resolution and accomplishment of the open requests. \\
IT Department should aim to higher percentage expecially related to high priority tasks.

\begin{table}[H]
    \centering
    \renewcommand{\arraystretch}{1.6}
    \setlength{\arrayrulewidth}{0.6pt}

    \begin{tabularx}{\textwidth}{|X|X|}
    \hline
    \multicolumn{2}{|c|}{\cellcolor{mcblue}\textcolor{white}{\textbf{Task Completed Since Last Report}}} \\ \hline

    \rowcolor{mclightgray}
    \textbf{Description} & \textbf{Unit} \\ \hline
    It indicates the team`s throughput and ability to reduce backlog. High percentage means high efficiency of the teams` abilities &
    Percentage. \\ \hline

    \rowcolor{mclightgray}
    \textbf{Formula} & \textbf{Frequency} \\ \hline
    \[
     \frac{\text{Tickets completed during last report}}{Tickets opened during last report} \times 100
    \]
    \textit{divided for each priority level if used as comparison within the previous KPI}
     &
    Montly (within the previous management report). \\ \hline

    \end{tabularx}

    \caption{Backlog - Task Completed Since Last Report}
    \label{tab:TaskCompletedLastReport}
\end{table}